% Teaching manuscript. Edit this TeX source; Markdown is compatibility navigation only.


\StudyChapter{CS5489-Tutorial04}{Tutorial 4 · Predicting Bike Sharing Usage}
\StudyTarget{CS5489-Tutorial04__cell-1}
\StudyTarget{CS5489-Tutorial04__notebook-c5900cbe}
\StudyTarget{CS5489-Tutorial04__tutorial-4-predicting-bike-sharing-usage}



\StudyHeading{1}{去年的规律，今年还管用吗？}
\StudyTarget{CS5489-Tutorial04___1}


\StudyRef{CS5489-Lecture04}{root}{Lecture 4} · \StudyRepoRef{CS5489/course-notes/README.md}{}{课程目录} · \StudyRef{Foundations-MathForML}{rank-pseudoinverse}{矩阵与回归基础}


老师给出两年的共享单车日记录，希望根据天气、日期类别等信息预测每日租车总量。这份学习笔记按原Tutorial4.ipynb的77个单元展开，单元从原文件第一个开始计数。数据直接从原ZIP读取，无需改动原文件。


前半练回归与特征解释，后半的关键问题是：第二年用户变多，旧模型可能系统性低估。给它补上第二年一季度的记录，怎样公平判断调整有没有帮助？


\textbf{基础入口（可跳过）：} 不熟数组形状先回\StudyRef{CS5489-Lecture01}{arrays}{Lecture 1}；不熟平方误差、正则项，先读\StudyRef{CS5489-Lecture04}{root}{Lecture 4}。先理解“输入的一行是一天，输出是当天总量”。


\begingroup
\small\setlength{\tabcolsep}{3pt}
\renewcommand{\arraystretch}{1.22}
\begin{longtable}{@{}>{\raggedright\arraybackslash}p{\dimexpr0.369111\linewidth-1.50000\tabcolsep\relax}>{\raggedright\arraybackslash}p{\dimexpr0.150658\linewidth-1.50000\tabcolsep\relax}>{\raggedright\arraybackslash}p{\dimexpr0.199531\linewidth-1.50000\tabcolsep\relax}>{\raggedright\arraybackslash}p{\dimexpr0.280700\linewidth-1.50000\tabcolsep\relax}@{}}
\toprule
原任务 & 优先级 & 难度／卡点 & 掌握要求与依据 \\
\midrule\endfirsthead
原任务 & 优先级 & 难度／卡点 & 掌握要求与依据 \\
\midrule\endhead
\bottomrule\endfoot
Tutorial4，第4–28个单元 数据、one-hot、年份划分 & 核心必会 & 中，特征与目标泄漏 & 能保留原2011/2012划分，解释删除计数列 \\
Tutorial4，第29–47个单元 基线、OLS、Ridge & 核心必会 & 中，CV与误差单位 & 能用训练内部验证选参数并读散点图 \\
Tutorial4，第48–58个单元 LASSO、OMP & 常规掌握 & 中，相关特征与系数含义 & 能比较选择结果，不把权重当因果 \\
Tutorial4，第59–64个单元 非线性回归 & 核心必会 & 中，公平模型比较 & 能固定划分、报告CV和测试各自用途 \\
Tutorial4，第65–76个单元 新年份与三种适应方法 & 核心必会 & 高，数据角色变化 & 能在同一剩余测试集比较四条流程 \\
\end{longtable}
\endgroup


优先级依据当前任务与Lecture4。代码与答案为参考实现；本讲义展示配套Notebook的已保存输出，复现时在Notebook中从头运行，独立尝试题在讲义中阅读。





\StudyHeading{1}{1. Setup｜先固定实验条件}
\StudyTarget{CS5489-Tutorial04__cell-2}
\StudyTarget{CS5489-Tutorial04__notebook-eabcab56}
\StudyTarget{CS5489-Tutorial04__1-setup}


原题未指定随机搜索范围。本参考使用固定种子42，2011年内的五折扩展窗口验证（TimeSeriesSplit），并让验证日期始终晚于对应训练日期。原全年2011训练、2012测试保持。模型参数只看2011验证误差选择；第二年成绩用于教学分析。


English: Hyperparameters are chosen within 2011 using forward chronological validation. The original year-based outer split is preserved. No model is selected using its 2012 test score.




\StudyTarget{CS5489-Tutorial04__cell-3}
\StudyTarget{CS5489-Tutorial04__notebook-e7c1c92b}
\StudyNotebookCell{CS5489/course-notes/Tutorial04.ipynb}{e7c1c92b}




\StudyTarget{CS5489-Tutorial04__cell-4}
\StudyTarget{CS5489-Tutorial04__notebook-f144725c}



\StudyTarget{CS5489-Tutorial04__2-load-and-inspectt428}




\StudyHeading{1}{2. Load and inspect｜计数目标与输入特征}
\StudyTarget{CS5489-Tutorial04__2-load-and-inspecttutorial4428}
\StudyTarget{CS5489-Tutorial04__2-load-and-inspect}


\textbf{任务 / Task：} 读取日数据，查看字段与分布，把season、mnth、weekday、weathersit做one-hot，预测cnt，按yr划分。 / Load daily records, inspect features, encode four categorical fields, predict cnt, and split by year.


ZIP同时有hour.csv和day.csv，本题用day.csv。\StudyInline{cnt=\allowbreak{}casual+registered}：若把后两列放入X，模型已经看到了答案的两个加数。预测当天总量时，还不知道当天两类人的完整数量，因此须从输入中删除这两列。


原说明把weekday写成1–7，实际数据是0–6；代码以数据为准。\StudyInline{instant}仅是记录号；原题删除日期列，但日期仍用于核验顺序。其他预测任务是否使用日期，应根据预测时可获得的信息决定。


来源：Tutorial4，第4–28个单元。




\StudyTarget{CS5489-Tutorial04__cell-5}
\StudyTarget{CS5489-Tutorial04__notebook-9dae4ba2}
\StudyNotebookCell{CS5489/course-notes/Tutorial04.ipynb}{9dae4ba2}




\StudyTarget{CS5489-Tutorial04__cell-6}
\StudyTarget{CS5489-Tutorial04__notebook-e087f529}
\StudyInline{yr}保留，因为后面合并新旧年份训练要用；只在2011年训练时，这一列全是0，模型没见过它从0变成1，因此学不到这个变化对应的需求差异。固定天气编码中值4在日数据里没有出现，对应列全0；这让列名稳定，不代表观察到了暴雪样本。天气值、连续特征已按数据说明缩放，不再从全部数据估计新的缩放参数。


\textbf{自测 / Check：} 保留registered会造成什么？ / What happens if registered remains an input?


答案 / Answer：它是当天总量的一部分，构成目标泄漏，测试成绩不再代表根据事前信息预测需求。 / It leaks part of the target and invalidates the intended advance-demand prediction.




\StudyTarget{CS5489-Tutorial04__cell-7}
\StudyTarget{CS5489-Tutorial04__notebook-64f6aae7}



\StudyTarget{CS5489-Tutorial04__3-a-baseline-and-interpretable-errorst2947}




\StudyHeading{1}{3. A baseline and interpretable errors｜常数基线与误差单位}
\StudyTarget{CS5489-Tutorial04__3-a-baseline-and-interpretable-errorstutorial42947}
\StudyTarget{CS5489-Tutorial04__3-a-baseline-and-interpretable-errors}


DummyRegressor只记训练目标的均值或中位数。它不使用天气等特征，能用来判断复杂模型是否确实增加了预测能力。MSE单位是“租车次数的平方”；RMSE开方后回到次数，MAE也是次数。原Tutorial4，第31个单元文字写MSE，Tutorial4，第32个单元实际代码计算RMSE；本笔记同时列出MSE、RMSE、MAE，避免把平方误差与开方后误差混用。


\textbf{任务 / Task：} 拟合基线、OLS、训练内选alpha的Ridge，并比较误差及预测散点。 / Fit dummy, OLS and CV-selected ridge models, then compare errors and scatter plots.


来源：Tutorial4，第29–47个单元。




\StudyTarget{CS5489-Tutorial04__cell-8}
\StudyTarget{CS5489-Tutorial04__notebook-52826fd9}
\StudyNotebookCell{CS5489/course-notes/Tutorial04.ipynb}{52826fd9}




\StudyTarget{CS5489-Tutorial04__cell-9}
\StudyTarget{CS5489-Tutorial04__notebook-695d8a3c}
\StudyNotebookCell{CS5489/course-notes/Tutorial04.ipynb}{695d8a3c}




\StudyTarget{CS5489-Tutorial04__cell-10}
\StudyTarget{CS5489-Tutorial04__notebook-67ff288d}
\begin{StudyBlock}{读图}{预测散点与残差的方向}
这里沿用老师Tutorial4，第32个单元的代码：横轴truth，纵轴prediction。因此低估在对角线\textbf{下方}，与Tutorial4，第66个单元一致。Tutorial4，第34个单元对基线的“x-value predictions”是文字简写，实际横轴仍是真值；读图以轴标签和代码为准。残差定义为真实值减预测值，正均值说明整体低估。


English: With truth on the horizontal axis and prediction on the vertical axis, underprediction lies below the diagonal. A positive mean residual (truth minus prediction) indicates underprediction.
\end{StudyBlock}

保存输出中的测试RMSE：OLS 2298.93、Ridge 2299.41，二者很接近，OLS略低。与均值基线2829.31相比两者均改善，但第二年的整体低估仍明显。 / In the saved results, OLS slightly outperforms ridge; both improve over the mean baseline, yet year-shift bias remains.




\StudyTarget{CS5489-Tutorial04__cell-11}
\StudyTarget{CS5489-Tutorial04__notebook-f11a88c0}



\StudyTarget{CS5489-Tutorial04__4-lasso-and-ompt4858}




\StudyHeading{1}{4. LASSO and OMP｜相关特征与系数解释}
\StudyTarget{CS5489-Tutorial04__4-lasso-and-omptutorial44858}
\StudyTarget{CS5489-Tutorial04__4-lasso-and-omp}


LASSO让一些权重恰为0；OMP逐次选择能改善残差的方向，再联合拟合已选特征。原题要求试10个特征，这里照做。\StudyInline{temp}与\StudyInline{atemp}、月份与季节可能相关，两种方法不必选同一批。


\textbf{任务 / Task：} 训练LASSO，按绝对系数排序并解释；用OMP选择10个特征，比较两组选择。 / Fit LASSO, interpret the largest absolute coefficients, and compare its selection with 10-feature OMP.


来源：Tutorial4，第48–58个单元。




\StudyTarget{CS5489-Tutorial04__cell-12}
\StudyTarget{CS5489-Tutorial04__notebook-b7a20806}
\StudyNotebookCell{CS5489/course-notes/Tutorial04.ipynb}{b7a20806}




\StudyTarget{CS5489-Tutorial04__cell-13}
\StudyTarget{CS5489-Tutorial04__notebook-0395c823}
比如温度权重为w，温度的归一化值增加0.1、其他输入不变时，预测变化为0.1w；不是增加1°C就变化w。One-hot权重的解释受基线、截距及同组其他类别影响，尤其完整one-hot加截距存在冗余。它描述给定其他输入时模型怎样改预测；实际骑行量还受其他因素影响。


English: Interpret a coefficient in the feature's encoded units, holding other inputs fixed. Correlated predictors can share or exchange weight; selection does not establish causation.


保存结果中LASSO选alpha=0.1，有28个非零系数；OMP按题设选10个。温度系数约3091.82、风速\ensuremath{-}1977.51、湿度\ensuremath{-}1281.77（均按原缩放单位）。选择不同并不自动说明谁算错，因为目标与稀疏约束不同。 / The fitted LASSO and ten-feature OMP use different sparsity objectives.




\StudyTarget{CS5489-Tutorial04__cell-14}
\StudyTarget{CS5489-Tutorial04__notebook-a1a2107d}



\StudyTarget{CS5489-Tutorial04__5-nonlinear-regressiont5964}




\StudyHeading{1}{5. Nonlinear regression｜让比较只改变模型}
\StudyTarget{CS5489-Tutorial04__5-nonlinear-regressiontutorial45964}
\StudyTarget{CS5489-Tutorial04__5-nonlinear-regression}


这里选原Lecture4中的RBF核岭回归与随机森林作两条参考路线。参数网格由本参考实现选定。每条路线用同一2011时间验证折，选定模型后才汇报2012。


\textbf{独立尝试 / Try：} 另加一种Lecture4模型，保持划分与误差单位，比较训练验证RMSE。 / Add another Lecture4 model while preserving splits and evaluation units.


来源：Tutorial4，第59–64个单元。




\StudyTarget{CS5489-Tutorial04__cell-15}
\StudyTarget{CS5489-Tutorial04__notebook-6e5179af}
\StudyNotebookCell{CS5489/course-notes/Tutorial04.ipynb}{6e5179af}




\StudyTarget{CS5489-Tutorial04__cell-16}
\StudyTarget{CS5489-Tutorial04__notebook-111ee21a}
训练得最好不等于跨年最好。树擅长分区间，RBF依靠与训练样本的相似性；面对整体需求变化，两者都可能失效。先按2011年的验证结果选模型，再观察它在2012年表现怎样；下面的结果展示了两者可能不同。


English: Compare the printed scores, and distinguish the CV-selected model from the model that happens to score best on the test table. Flexible functions alone do not solve temporal distribution shift.


2011时间验证选出LASSO，故后续适应实验固定它。事后2012表中随机森林RMSE最低（2205.32），这一事后排名不用于改变后续已固定的方法。非线性候选没有赢得本次训练内选择。 / LASSO wins the training-only selection; the test-table winner is not used to revise that choice.




\StudyTarget{CS5489-Tutorial04__cell-17}
\StudyTarget{CS5489-Tutorial04__notebook-ea2c8c13}



\StudyTarget{CS5489-Tutorial04__6-test-set-biast6576}




\StudyHeading{1}{6. Test set bias｜先把新增数据的角色划清楚}
\StudyTarget{CS5489-Tutorial04__6-test-set-biastutorial46576}
\StudyTarget{CS5489-Tutorial04__6-test-set-bias}


画总量随日期变化，再模拟老师给定的情境：2012年1–3月已经结束，其标签现在可用；把它当新训练数据，4–12月仍只作评估。不能一边用一季度调模型，一边把它算进最终测试成绩。


筛选第二年一季度时，参考实现直接使用原表的mnth月份字段，取得与课堂方案相同的记录：91天新训练，275天剩余测试。特征仍按同一列顺序。来源注：Tutorial4，第69个单元按第11–13列取月份；依赖列位置的写法需核对编码后的列顺序。


\textbf{任务 / Task：} 固定一种已选回归方法，比较旧模型、新数据单独重训、合并重训、学习旧模型残差；全部在同一个4–12月测试集评估。 / Compare the original model, new-only retraining, pooled retraining, and residual adaptation on the same April–December holdout.


来源：Tutorial4，第65–76个单元。




\StudyTarget{CS5489-Tutorial04__cell-18}
\StudyTarget{CS5489-Tutorial04__notebook-f0d1eeb1}
\StudyNotebookCell{CS5489/course-notes/Tutorial04.ipynb}{f0d1eeb1}




\StudyTarget{CS5489-Tutorial04__cell-19}
\StudyTarget{CS5489-Tutorial04__notebook-cb7ee79e}
\begin{StudyBlock}{方法条件}{残差修正的符号与时间范围}
残差修正保留旧模型，由新模型预测旧模型尚未解释的差值。约定\(r=y-f_{old}(x)\)时，加上\(\hat r\)；若把残差反过来定义为\(f_{old}(x)-y\)，修正时必须减。原题文字没有把这个符号约定讲透，实施时不能前后混用。


三种方法都不保证最好。新数据只有冬末春初，单独训练可能看不到夏季；合并数据覆盖更广，但旧年份占比大；残差法在旧预测上加修正，仍可能改变季节曲线，并依赖“差值规律能迁移到后九个月”的假设。比较时结合上表误差和这些假设解释结果。
\end{StudyBlock}

\textbf{迁移题 / Transfer：} 若现在只到2012年2月底，可以拿3月数据选参数吗？为什么？ / If it is only the end of February 2012, may March observations be used for tuning an operational predictor?



\begin{StudySolution}{答案 / Answer}


不能，3月结果尚不可得；应把可用日期以前的数据进一步按时间拆分。离线已下载全年数据不代表当时可见。 / No. March outcomes are not yet available. Validation must reflect the information available at the prediction date.


\end{StudySolution}


在保存结果的同一275天测试集上，合并重训的RMSE最低，为1042.12。残差修正仍整体高估，平均真值减预测为\ensuremath{-}1362.37，提示早春的偏差规律未充分代表其余季节。 / In the saved results for the same 275 held-out days, pooled retraining has the lowest RMSE (1042.12). Residual adaptation still overpredicts, with a mean truth-minus-prediction residual of \ensuremath{-}1362.37.





\StudyHeading{1}{7. Return to the lecture｜把结果变成能解释的结论}
\StudyTarget{CS5489-Tutorial04__cell-20}
\StudyTarget{CS5489-Tutorial04__notebook-8fb0865c}
\StudyTarget{CS5489-Tutorial04__7-return-to-the-lecture}


英文自测 / English self-test：Why were the count components removed? Why split by year? Why does a constant year feature fail to teach the next year's offset? Why must all adaptation methods share the same holdout?


中文题意：为什么删去两项计数组成？为什么按年份划分？训练中恒定的年份特征为何学不到下一年偏移？为什么适应方法必须使用同一测试集？


关键答案 / Key answers：避免目标泄漏；评估向未来迁移；训练中无变化的特征无法识别变化效应；同一测试集才有可比较的误差。 / Prevent target leakage; evaluate future transfer; a constant feature cannot identify a changing effect; compare errors on identical held-out records.


首轮复习沿用\href{https://crazyshout.github.io/micro-course/cards.html\#CS5489-M187}{回归 M187}、\href{https://crazyshout.github.io/micro-course/cards.html\#CS5489-M194}{Ridge M194}、\href{https://crazyshout.github.io/micro-course/cards.html\#CS5489-M196}{LASSO M196}，再回\href{https://crazyshout.github.io/micro-course/?lesson=ml19}{ML19}串起目标、约束与验证。


\begingroup
\small\setlength{\tabcolsep}{3pt}
\renewcommand{\arraystretch}{1.22}
\begin{longtable}{@{}>{\raggedright\arraybackslash}p{\dimexpr0.714244\linewidth-1.00000\tabcolsep\relax}>{\raggedright\arraybackslash}p{\dimexpr0.285756\linewidth-1.00000\tabcolsep\relax}@{}}
\toprule
原课单元 & 本笔记位置 \\
\midrule\endfirsthead
原课单元 & 本笔记位置 \\
\midrule\endhead
\bottomrule\endfoot
Tutorial4，第1–3个单元 姓名/任务/初始化 & 开头与§1，个人身份不填入公共学习稿 \\
Tutorial4，第4–28个单元 原数据、分布、统计、编码、提取与年份划分 & §2 \\
Tutorial4，第29–35个单元 Dummy与误差/散点 & §3 \\
Tutorial4，第36–47个单元 OLS/Ridge与解释 & §3 \\
Tutorial4，第48–58个单元 LASSO/OMP与系数 & §4 \\
Tutorial4，第59–64个单元 非线性模型 & §5 \\
Tutorial4，第65–76个单元 趋势及三种再训练 & §6 \\
Tutorial4，第77个单元 空单元 & 无教学内容 \\
\end{longtable}
\endgroup


结果由以上代码产生，采用§1的固定设置。




\StudyTarget{CS5489-Tutorial04__cell-21}
\StudyTarget{CS5489-Tutorial04__notebook-b976a989}
\StudyNotebookCell{CS5489/course-notes/Tutorial04.ipynb}{b976a989}
